\documentclass[10pt,a4paper]{amsart}
\usepackage[margin=19mm]{geometry}
\usepackage{amsmath,amssymb,mathtools}
\usepackage[scr=rsfs,cal=euler]{mathalfa}
\usepackage{xfrac}
\usepackage[unicode,hidelinks,bookmarks=false]{hyperref}
\newcommand{\ie}{i.e.\ }
\newcommand{\cf}{cf.\ }
\newcommand{\resp}{respectively\ }
\newcommand{\Subsection}{\subsection}
\providecommand{\nobreakdash}{\nobreak\hbox{-}}
\setlength{\emergencystretch}{2em}
\multlinegap0pt
\usepackage{bm}
\usepackage{bbm}
\usepackage{dsfont}
\usepackage{xspace}
\usepackage{cancel}
\usepackage{mathtools}
\usepackage{amsthm}
\makeatletter
\@ifundefined{assumption}{\newtheorem{assumption}{Assumption}}{}
\@ifundefined{theorem}{\newtheorem{theorem}{Theorem}}{}
\makeatother
\makeatletter
\newcommand{\Lip}{\operatorname{Lip}}
\newcommand{\RR}{\mathbb{R}}
\expandafter\ifx\csname method\endcsname\relax
\newenvironment{method}[1]{%
	\refstepcounter{method}%
	\par\medskip
	\noindent\textbf{Method~\themethod: #1.}\ }%
\fi
\newcommand{\supp}{\operatorname{supp}}
\makeatother
\title[A Unified DeepONet Framework for Logarithmically Stable Infi]{A Unified DeepONet Framework for Logarithmically Stable Infinite-Dimensional Inverse Problems}
\author{Wen-Jie Wu, Tiexiang Li, Wen-Wei Lin}
\date{Source: arXiv:2606.07122v2 (CC BY 4.0)}
\begin{document}
\maketitle

\noindent\emph{Source and licence.} A Unified DeepONet Framework for Logarithmically Stable Infinite-Dimensional Inverse Problems. Version arXiv:2606.07122v2 (CC BY 4.0), distributed under the Creative Commons Attribution 4.0 International licence, as recorded in the arXiv licence metadata for this exact version and shown on the abstract page https://arxiv.org/abs/2606.07122v2. Full rights notice: licenses/arxiv-2606.07122-CC-BY-4.0.txt.\par\smallskip

\noindent\textit{Editorial scope.} Counted unit: the Theorem whose source label is \texttt{thm:error-decomposition} in the article (source0.tex lines 1162--1206, source label \texttt{thm:error-decomposition}), delivered together with the article's own complete original proof of it (source0.tex lines 1208--1380, one \texttt{proof} environment of 173 lines). No mathematics is written, repaired or re-derived: statement and proof are the article's own text line for line. Required context retained uncounted: ass:stability (lines 263--276); thm:abstract-approximation (lines 410--449). Every definition, display and label of those lines is retained unchanged. Omitted from this package and not redistributed: 1--262, 277--409, 450--1161, 1207--1207, 1381--2816. Nothing inside a retained excerpt is omitted or altered: every retained body is delivered exactly as the article's lines are written, so no omission receipt of any kind accompanies this package. Citations: the retained excerpts carry no citation command at all, so no bibliography entry of the article is transcribed and no reference is redistributed. Reference adaptation: none was needed and none was made; every reference inside the delivered unit resolves inside it, so no unresolved reference or citation is produced. Printed slips kept literal: no printed word, symbol, hypothesis or number of the counted statement or of its proof was altered. Layout and class: the article's own class is \texttt{article} with packages including fontenc, inputenc, amsmath, amssymb, amsthm, geometry, mathrsfs, xcolor, microtype, tikz, graphicx, booktabs, hyperref, float; the wrapper is the lane's pinned \texttt{amsart} editorial wrapper at 10pt on A4, which restates the theorem-like environments the retained text needs and adds the article's own macro definitions that the retained text needs (\texttt{\textbackslash Lip}, \texttt{\textbackslash RR}, \texttt{\textbackslash supp}), with extra wrapper packages (bm, bbm, dsfont, xspace, cancel, mathtools). The delivered numbering is the unit's own. Assets: no retained span contains a figure environment, a graphics-inclusion command, a PDF-page inclusion, a table or a TikZ picture; no image file is copied into this package, the package declares no asset dependency and no image file is redistributed with it. Licence: version arXiv:2606.07122v2 of this article is distributed under the Creative Commons Attribution 4.0 International licence, as recorded in the arXiv licence metadata for this exact version and shown on the abstract page https://arxiv.org/abs/2606.07122v2. A full rights notice is delivered at licenses/arxiv-2606.07122-CC-BY-4.0.txt. Characters: every retained excerpt is pure ASCII, so the delivered engine, XeLaTeX with its default Latin Modern text font, reports no missing character.
	\begin{assumption}[Logarithmic stability]\label{ass:stability}
		Assume that the inverse map
		\[
		G:D_G\to Y
		\]
		satisfies the following logarithmic stability estimate: there exist constants
		$C_{\mathrm{stab}}>0$, $\beta\in(0,1)$, and $\tau\in(0,1)$ such that
		\[
		\|G(x_1)-G(x_2)\|_Y
		\le
		C_{\mathrm{stab}}\,|\log \|x_1-x_2\|_X|^{-\beta},
		\]
		for all $x_1,x_2\in D_G$ satisfying $\|x_1-x_2\|_X\le \tau$. Here and below, we use the continuous extension convention $|\log 0|^{-\alpha}:=0$ for every $\alpha>0$.
	\end{assumption}

	\begin{theorem}[Approximation error under logarithmic stability]\label{thm:abstract-approximation}
		Fix $m,p\in\mathbb N$. Suppose Assumption~\ref{ass:stability} holds. Let $\mathcal D:\mathbb R^m\to X$ and $\mathcal P:Y\to\mathbb R^p$ be Lipschitz continuous maps, and let $K:=\supp(\mathcal E_\#\mu)\subset \mathbb R^m$ be compact. Assume that $\mathcal D(K)\subset D_G$, and define $H:=\mathcal P\circ G\circ\mathcal D:K\to \mathbb R^p$.
		Then for every $\eta\in(0,1)$ there exists a ReLU neural network $\mathcal A_\eta:\mathbb R^m\to \mathbb R^p$
		such that
		\begin{equation*}
			\mathscr{E}_{\mathcal{A}}(\mathcal A_\eta)
			:=\|\mathcal A_\eta - H\|_{L^{2}(K,\,\mathcal{E}_{\#}\mu)}
			\le\eta,
		\end{equation*}
		and
		\begin{equation*}
			\mathrm{size}(\mathcal A_\eta)
			\le c_0(p+1)\exp \bigl(c_1m\eta^{-1/\beta}\bigr),
			\qquad
			\mathrm{depth}(\mathcal A_\eta)
			\le c_2+c_3m\eta^{-1/\beta}.
		\end{equation*}
		The constants $c_0,c_1,c_2,c_3>0$ may depend on the prescribed ranks $m,p$ and on the fixed data in the assumptions, but are independent of $\eta$ and $N$.
		For $N\in\mathbb N$, define
		\begin{equation*}
			\mathscr E_{\mathcal A}(N)
			:=
			\inf_{\substack{\mathcal A:\mathbb R^m\to\mathbb R^p\ \text{is a ReLU network}\\
			\mathrm{size}(\mathcal A)\le N}}
			\|\mathcal A-H\|_{L^2(K,\mathcal E_\#\mu)}.
		\end{equation*}
		Set
		\begin{equation*}
			N_0:=c_0(p+1)\exp(c_1m).
		\end{equation*}
		Then, for every $N> N_0$,
		\begin{equation*}
			\mathscr E_{\mathcal A}(N)
			\le
			\left[
			\frac{1}{c_1m}
			\log\!\left(\frac{N}{c_0(p+1)}\right)
			\right]^{-\beta}.
		\end{equation*}
	\end{theorem}

	\begin{theorem}[DeepONet error estimate]\label{thm:error-decomposition}
		Fix $m,p\in\mathbb N$, and let $c_0,c_1,c_2,c_3,N_0$ be as in Theorem~\ref{thm:abstract-approximation}. Assume the logarithmic stability estimate in Assumption~\ref{ass:stability}. Let $\mu$ be a Borel probability measure on $X$ such that $\mu(D_G)=1$, and let
		\begin{equation*}
			K:=\operatorname{supp}(\mathcal E_\#\mu)\subset \mathbb R^m
		\end{equation*}
		be compact. Assume that $\mathcal D(K)\subset D_G$, that $\mathcal D:\mathbb R^m\to X$ and $\mathcal P:Y\to\mathbb R^p$ are Lipschitz continuous, and that
		\begin{equation*}
			\varepsilon_m
			:=
			\operatorname*{ess\,sup}_{x\in X}
			\|\mathcal D\circ\mathcal E(x)-x\|_X
			\le \min\{\tau,e^{-1}\}.
		\end{equation*}
		The essential supremum is taken with respect to $\mu$, and hence $\mathscr E_{\mathcal E}\le\varepsilon_m$ since $\mu$ is a probability measure. Assume further that $\mathcal R$ is Lipschitz continuous. In addition, suppose that, for some prescribed reconstruction rate $\varepsilon_R(p)$, the chosen reconstruction mechanism satisfies $\mathscr E_{\mathcal R} \le C_2 \varepsilon_R(p)$ for some constant $C_2>0$ independent of $p$.
		
		Then, for every $N>N_0$, there exists a ReLU approximator
		$\mathcal A_N:\mathbb R^m\to\mathbb R^p$, with
		$\mathrm{size}(\mathcal A_N)\le N$, such that
		\begin{equation}\label{total-error}
			\mathscr E
			\le
			\Lip(\mathcal R)
			\left[
			\frac{1}{c_1m}
			\log\!\left(\frac{N}{c_0(p+1)}\right)
			\right]^{-\beta}
			+
			\Lip(\mathcal R\circ\mathcal P)\,C_{\mathrm{stab}}|\log\varepsilon_m|^{-\beta}
			+
			C_2\varepsilon_R(p).
		\end{equation}
		Moreover,
		%\begin{enumerate}
			%\item $\mathscr E_{\mathcal E}\le\varepsilon_m$;
			%\item $\mathscr E_{\mathcal R}\le C_2\varepsilon_R(p)$;
			for every $\eta\in(0,1)$, the construction provides an approximator satisfying
			\begin{equation*}
				\mathrm{size}(\mathcal A_\eta)
				\le c_0(p+1)\exp \bigl(c_1m\eta^{-1/\beta}\bigr),
				\qquad
				\mathrm{depth}(\mathcal A_\eta)
				\le c_2+c_3m\eta^{-1/\beta}.
			\end{equation*}
		%\end{enumerate}
	\end{theorem}

	\begin{proof}
		We decompose the total error into three parts:
		\begin{equation*}
			\begin{aligned}
				\mathscr{N} - \mathscr{G}
				&= \mathcal{R}\circ\mathcal{A}\circ\mathcal{E} - \mathscr{G} \\
				&= [\mathcal{R}\circ\mathcal{A}\circ\mathcal{E} - \mathcal{R}\circ\mathcal{P}\circ \mathscr{G}]
				+[\mathcal{R}\circ\mathcal{P}\circ \mathscr{G}-\mathscr{G}] \\
				&= [\mathcal{R}\circ\mathcal{A}\circ\mathcal{E} - \mathcal{R}\circ\mathcal{P}\circ \mathscr{G}\circ\mathcal{D}\circ\mathcal{E}] \\
				&\quad + [\mathcal{R}\circ\mathcal{P}\circ \mathscr{G}\circ\mathcal{D}\circ\mathcal{E} - \mathcal{R}\circ\mathcal{P}\circ \mathscr{G}] \\
				&\quad + [\mathcal{R}\circ\mathcal{P}\circ \mathscr{G} - \mathscr{G}] \\
				&=: T_1+T_2+T_3.
			\end{aligned}
		\end{equation*}
		
		We can now estimate the norms of these three terms as follows:
		\begin{equation*}
			\begin{aligned}
				\|T_1\|_{L^2(\mu)}
				&=
				\left(
				\int_X
				\|\mathcal{R}\circ\mathcal{A}\circ\mathcal{E}
				-
				\mathcal{R}\circ\mathcal{P}\circ \mathscr{G}\circ\mathcal{D}\circ\mathcal{E}\|_{Y}^{2}
				\,d\mu
				\right)^{1/2}\\
				&\leq
				\Lip(\mathcal{R})
				\left(
				\int_X
				\|\mathcal{A}\circ\mathcal{E}
				-
				\mathcal{P}\circ \mathscr{G}\circ\mathcal{D}\circ\mathcal{E}\|_{\ell^2(\RR^p)}^{2}
				\,d\mu
				\right)^{1/2}\\
				&=
				\Lip(\mathcal{R})
				\left(
				\int_{\RR^m}
				\|
				\mathcal{A}
				-
				\mathcal{P}\circ \mathscr{G}\circ\mathcal{D}
				\|_{\ell^2(\RR^p)}^{2}
				\,d(\mathcal{E}_{\#}\mu)
				\right)^{1/2}\\
				&=
				\Lip(\mathcal{R})\cdot \mathscr{E}_{\mathcal{A}}.
			\end{aligned}
		\end{equation*}
		
		For the second term, we get
		\begin{equation*}
			\begin{aligned}
				\|T_2\|_{L^2(\mu)}
				&=
				\left(
				\int_X
				\|
				\mathcal{R}\circ\mathcal{P}\circ \mathscr{G}\circ\mathcal{D}\circ\mathcal{E}
				-
				\mathcal{R}\circ\mathcal{P}\circ \mathscr{G}
				\|_{Y}^{2}
				\,d\mu
				\right)^{1/2}\\
				&\leq
				\Lip(\mathcal{R}\circ\mathcal{P})
				\left(
				\int_X
				\|
				\mathscr{G}\circ\mathcal{D}\circ\mathcal{E}
				-
				\mathscr{G}
				\|_{Y}^{2}
				\,d\mu
				\right)^{1/2}.
			\end{aligned}
		\end{equation*}
		
		For $\mu$-a.e. $x$, both $x$ and $\mathcal D\circ\mathcal E(x)$ belong to $D_G$. Applying the logarithmic stability estimate,
		\begin{equation*}
			\|\mathscr{G}\circ\mathcal{D}\circ\mathcal{E}(x)-\mathscr{G}(x)\|_Y
			\leq
			C_{\mathrm{stab}}\,\omega(\|\mathcal{D}\circ\mathcal{E}(x)-x\|_X),
		\end{equation*}
		where $\omega(t)=|\log t|^{-\beta}$.
		Let
		\begin{equation*}
			\delta(x)=\|\mathcal{D}\circ\mathcal{E}(x)-x\|_X,
		\end{equation*} 
		\begin{equation*}
			\varepsilon_m:=\operatorname*{ess\,sup}_{x\in X}\|\mathcal{D}\circ\mathcal{E}(x)-x\|_X,
		\end{equation*}
		where the essential supremum is taken with respect to $\mu$.
		By definition of $\varepsilon_m$, $\delta(x)\le \varepsilon_m$ for $\mu$-a.e. $x$. Since $\varepsilon_m\le \min\{\tau,e^{-1}\}$, the logarithmic stability estimate applies on the relevant range and $\omega$ is monotone there. Hence
		\begin{equation*}
			\omega(\delta(x))
			\leq
			\omega(\varepsilon_m)
			=
			|\log \varepsilon_m|^{-\beta}.
		\end{equation*}
		Therefore,
		\begin{equation*}
			\left(
			\int_X (\omega(\delta(x)))^2\,d\mu(x)
			\right)^{1/2}
			\leq
			\omega(\varepsilon_m)
			=
			|\log \varepsilon_m|^{-\beta}.
		\end{equation*}
		We obtain
		\begin{equation*}
			\|T_2\|_{L^2(\mu)}
			\leq
			\Lip(\mathcal{R}\circ\mathcal{P})\cdot C_{\mathrm{stab}}\cdot |\log \varepsilon_m|^{-\beta}.
		\end{equation*}
		
		For the third term, we obtain
		\begin{equation*}
			\begin{aligned}
				\|T_3\|_{L^2(\mu)}
				&=
				\left(
				\int_X
				\|
				\mathcal{R}\circ\mathcal{P}\circ \mathscr{G}
				-
				\mathscr{G}
				\|_{Y}^{2}
				\,d\mu
				\right)^{1/2}\\
				&=
				\left(
				\int_Y
				\|(\mathcal R\circ\mathcal P)(y)-y\|_Y^2
				\,d(\mathscr G_\#\mu)(y)
				\right)^{1/2}\\
				&=
				\mathscr E_{\mathcal R}.
			\end{aligned}
		\end{equation*}
		By the assumed reconstruction error estimate, we have
		\begin{equation*}
			\|T_3\|_{L^2(\mu)}
			\leq
			C_2\,\varepsilon_R(p).
		\end{equation*}
		
		Combining these estimates yields
		\begin{equation*}
			\mathscr{E}
			\leq
			\Lip(\mathcal{R})\cdot \mathscr{E}_{\mathcal{A}}
			+
			\Lip(\mathcal{R}\circ\mathcal{P})\cdot C_{\mathrm{stab}}\cdot |\log \varepsilon_m|^{-\beta}
			+
			C_2\,\varepsilon_R(p).
		\end{equation*}
		
		Finally, Theorem~\ref{thm:abstract-approximation} provides an approximator $\mathcal A_N$ of size at most $N$ satisfying
		\begin{equation*}
			\mathscr E_{\mathcal A}
			\le
			\left[
			\frac{1}{c_1m}
			\log\!\left(\frac{N}{c_0(p+1)}\right)
			\right]^{-\beta}.
		\end{equation*}
		Substitution into the preceding decomposition gives \eqref{total-error}.
	\end{proof}

\end{document}
